% ================================================================
% settings.tex – Pakete, Layout, Kopf-/Fußzeile (Berufsschul-Skript)
% (Titel, Lernfeld usw. stehen in hauptdokument.tex)
% ================================================================

%%%%%%%%%%%%%%%%%%%%%%%%%%%%%% Pakete %%%%%%%%%%%%%%%%%%%%%%%%%%%%%%%%%%%
\usepackage{ngerman}
\usepackage[utf8]{inputenc}
\usepackage[T1]{fontenc}
\usepackage{graphicx}
\usepackage{lastpage}
\usepackage{float}
\usepackage{caption}
\usepackage[]{todonotes} % [disable] blendet alle \todo aus
\usepackage{dialogue}
\usepackage{multicol}
\usepackage{pdfpages}

\usepackage[dvipsnames]{xcolor}
\usepackage{amsmath}
\usepackage{amssymb}
\usepackage{gensymb}
\usepackage{cancel}

\usepackage{enumitem}
\usepackage{acronym}
\usepackage{listings}

% Tabellen
\usepackage{setspace}  % Zeilenhöhe
\usepackage{colortbl}  % Hintergrundfarben in Tabellen
\usepackage{tabularx}
\usepackage{array}     % z. B. \newcolumntype

% Schrift: serifenlos (für Serifen die renewcommand-Zeile auskommentieren)
\usepackage{lmodern}
\renewcommand{\familydefault}{\sfdefault}

%%%%%%%%%%%%%%%%%%%%%%%%%%%%%% Eigene Definitionen %%%%%%%%%%%%%%%%%%%%%%
% Farben, Boxen, Listings, Aufgaben, Pfade (\imgpath, \filepath, \lstpath, \sharedpath)
\newcommand{\basispath}{/Users/florian/workspace/latex/vorlagen/latex-basis}
\input{\basispath/basis}

%%%%%%%%%%%%%%%%%%%%%%%%%%%%%% Hyperref (möglichst spät laden) %%%%%%%%%%%
\usepackage{hyperref}
\hypersetup{
    colorlinks=true,
    linkcolor=DRSblue,
    filecolor=DRSyellow,
    citecolor=DRSred,
    urlcolor=DRSblue,
    pdftitle={\scripttitel},
    pdfpagemode=FullScreen,
}

%%%%%%%%%%%%%%%%%%%%%%%%%%%%%% Page Layout %%%%%%%%%%%%%%%%%%%%%%%%%%%%%%
\usepackage{geometry}
\geometry{
  a4paper,
  top=35mm,
  left=25mm,
  right=25mm,
  bottom=30mm,
  headheight=53pt  % Platz für die dreizeilige Kopfzeile mit Logo
}

%%%%%%%%%%%%%%%%%%%%%%%%%%%%%% Kopf- und Fußzeilen %%%%%%%%%%%%%%%%%%%%%%
\usepackage[headsepline, automark]{scrlayer-scrpage}
\automark{section}
\setheadsepline{1pt}
\pagestyle{scrheadings}

\ohead{\includegraphics[height=1.7cm]{\boxenimages/drslogo}}
\ihead{Lernfeld \lernfeldnr : \lernfeldtext \\
  \bildungsgang \\
  \scripttitel
}
\chead{}

% \leftmark = aktueller Teil, automatisch aus \section (\automark{section})
\ifoot{Autor: \lehrkraft \\
  Stand: \today \\
  \leftmark}
\cfoot{Seite \thepage\ von \pageref{LastPage}}
%\ofoot{\thepage}

%%%%%%%%%%%%%%%%%%%%%%%%%%%%%% Einstellungen %%%%%%%%%%%%%%%%%%%%%%%%%%%%
% Musterlösung ein-/ausblenden: \ifloesung ... \fi
% (einschalten mit \loesungtrue in hauptdokument.tex)
\newif\ifloesung

% Kein Absatzeinzug
\setlength{\parindent}{0pt}
